\documentclass[10pt]{article}
\usepackage[margin=0.72in]{geometry}
\usepackage[T1]{fontenc}
\usepackage{lmodern}
\usepackage{microtype}
\usepackage{booktabs}
\usepackage{array}
\usepackage{amsmath}
\usepackage{amssymb}
\usepackage{amsthm}
\usepackage{hyperref}

\newtheorem{theorem}{Theorem}
\newtheorem{lemma}{Lemma}
\newtheorem{definition}{Definition}

\title{PromptABI: Static Verification for the Discrete Interface Layer of LLM Systems}
\author{PromptABI Contributors}
\date{}

\begin{document}
\maketitle

\begin{abstract}
Modern LLM systems fail in places that are not neural: chat templates splice
untrusted text into role delimiters, tokenizer metadata changes the strings that
grammars actually see, stop sequences fire inside tool arguments, provider
adapters reinterpret structured outputs, and framework truncation silently drops
the evidence a prompt was meant to preserve. These failures are discrete,
versioned, and usually present before a single token is generated, yet they are
typically discovered through inference-time tests or production incidents.

PromptABI treats this boundary as an application binary interface for prompts.
It gives templates, tokenizers, stop policies, tool schemas, provider envelopes,
retrieval budgets, training manifests, and evaluation harnesses a common typed
semantics. It then runs bounded symbolic execution, finite-language products,
differential replay, and SMT-backed finite contracts to prove one of four useful
facts: the interface is safe in the supported fragment, a concrete counterexample
exists, the checker must abstain, or a versioned artifact has drifted from a
previously verified contract.

The result is a CPU-only verifier that can be run in CI, before fine-tuning,
before publishing benchmark scores, or before migrating providers. Its evidence
is not a checklist of repository files: diagnostics carry executable witnesses,
source spans, token paths, parser states, solver models, replay hashes, and
privacy-preserving fingerprints. The reported bug-finding evaluation uses public
bug replays that traverse exact pinned production-code excerpts.
\end{abstract}

\section{Introduction}

The software around an LLM has become a compiler pipeline. A user message is
rendered by a chat template, transformed by tokenizer normalization and special
token rules, packed into a context window, constrained by a grammar or JSON
schema, streamed through stop policies, decoded by a provider adapter, and
possibly reused as training or evaluation data. The model may be probabilistic,
but this surrounding pipeline is largely deterministic. When it is wrong, the
failure is often predictable without running the model.

PromptABI is built on that observation. It verifies the discrete interface layer
of LLM systems before inference. The name is intentional: like an ABI for native
code, the prompt interface defines what representations cross subsystem
boundaries and which invariants downstream components assume. If a prompt
template says that user content is data, but renders raw role delimiters inside
that content, the ABI is broken. If a tool parser assumes that
\texttt{</tool\_call>} is structural, but a valid string argument may contain
that sequence, the ABI is broken. If an evaluation harness scores a model after
dropping the system prompt under a framework's packing policy, the ABI is
broken.

This paper presents PromptABI as a programming-languages tool for LLM
applications. The focus is not model quality, prompt style, or policy
moderation. It is the statically checkable semantics of the artifacts that
connect a model to the rest of a program.

\paragraph{Contributions.}
PromptABI contributes:
\begin{enumerate}
  \item a typed interface model for prompt-system artifacts that are otherwise
  checked independently or only by example;
  \item bounded semantics for chat-template rendering, tokenization, stop
  reachability, structured outputs, provider migration, context budgets,
  training masks, and evaluation harnesses;
  \item proof-producing diagnostics that distinguish verified safety,
  concrete counterexamples, and honest abstention;
  \item a deterministic CLI and embedding API designed for CI and offline use;
  \item an evaluation suite whose reported bug-finding results come from public
  production code in the wild, not synthetic simplifications of those bugs.
\end{enumerate}

\section{The Need}

LLM applications are commonly tested end-to-end: run a model, inspect a sample,
and patch the prompt when behavior looks wrong. That style misses interface
bugs for three reasons.

\paragraph{The bug may be rare under sampling but inevitable in the interface.}
If a stop string is allowed inside a valid tool argument, generation may or may
not produce it in a test run. The interface, however, admits the bad string
regardless of the model distribution.

\paragraph{The responsible component is often not the model.}
Many incidents are caused by a tokenizer revision, chat-template update,
provider adapter, streaming parser, framework truncation policy, or evaluation
wrapper. These artifacts are versioned software inputs and can be checked like
software inputs.

\paragraph{A passing run is not a proof.}
An inference test demonstrates one execution. A verifier can instead enumerate
the bounded states of the interface, produce a minimal witness, and record why a
claim is valid only in a supported fragment.

PromptABI addresses these gaps by moving checks left: the same artifacts that
would be passed to inference are loaded, normalized, and checked before any
weights or provider credentials are needed.

\section{A Prompt ABI}

\begin{definition}[Prompt interface artifact]
A prompt interface artifact is a discrete, versioned object that influences the
string or structured value seen by an LLM subsystem. Examples include chat
templates, tokenizer configs, stop policies, JSON schemas, tool definitions,
provider request/response envelopes, truncation policies, prompt segments,
training manifests, and evaluation manifests.
\end{definition}

\begin{definition}[Prompt ABI]
A prompt ABI is the finite contract induced by a set of prompt interface
artifacts: which roles may appear, which strings are data versus control, which
tokens are reserved, which regions must survive packing, which grammars can be
accepted, which stops may fire, and which structured fields downstream parsers
expect.
\end{definition}

PromptABI compiles each artifact into small semantic views. A chat template
becomes a bounded set of symbolic render paths and role regions. A tokenizer
becomes an encode/normalize/decode relation over finite witnesses. A JSON
schema becomes a grammar fragment and representative values. A provider adapter
becomes a contract over fields, tool-call encodings, streaming chunks, stops,
and structured-output modes. A training or evaluation manifest becomes a set of
finite obligations about roles, masks, preserved context, answer schemas, and
private fields.

The tool deliberately does not claim to model arbitrary Python, arbitrary
Jinja, or arbitrary parser code. Its supported fragments are explicit. Outside
them, PromptABI reports abstention instead of returning a false proof.

\section{How PromptABI Checks Interfaces}

PromptABI's checkers share a common pattern:
\begin{enumerate}
  \item load typed artifacts with source spans, hashes, and version metadata;
  \item compile them to bounded semantic objects;
  \item run a product, symbolic execution, differential replay, or SMT query;
  \item emit either a proof summary, a concrete witness, or an abstention.
\end{enumerate}

\subsection{Role Boundaries}

The role-boundary checker symbolically executes supported chat-template
fragments. It partitions each rendered path into structural regions such as
system, user, assistant, tool, developer, and function. It then asks whether an
attacker-controlled field can render a structural marker: role headers, special
tokens, assistant prefixes, tool delimiters, markdown fences, or provider
sentinels.

For example, a template that renders
\[
\texttt{'<|user|>\textbackslash n' + message['content'] + '<|end|>\textbackslash n'}
\]
without escaping permits user content to contain \texttt{<|system|>} or
\texttt{<|end|>} as literal control text. PromptABI reports the input field,
the marker, the rendered excerpt, the tokenized representation, and the exact
offset of the forged boundary.

\subsection{Stops and Structured Outputs}

Stop policies are safe only relative to the language of outputs they interrupt.
PromptABI builds bounded structured-output regions for JSON strings, tool
arguments, XML-like tool calls, markdown fences, and provider tool-call
envelopes. It then checks whether a configured stop sequence can occur before a
valid structured output is complete.

The witness records the stop, the valid full output, the prefix returned after
truncation, the parser state at the firing point, and the malformed or
prematurely accepted structure a downstream parser would receive.

\subsection{Tokenizers and Grammars}

Tokenizer bugs often arise from mismatches between text-level constraints and
token-level behavior. PromptABI can compare tokenizer revisions, detect drift in
special-token IDs and normalization, and form bounded products between a
tokenizer relation and a grammar. This exposes empty products, ambiguous token
paths, decoded-text conflicts, whitespace aliases, Unicode normalization
surprises, and added-token interactions.

\subsection{Budgets, Training, and Evaluation}

Context windows are treated as finite packing problems. PromptABI combines
declared segments, tokenizer-derived or declared token counts, output reserves,
tool overhead, generation prompts, and framework truncation policies to prove
whether required regions survive. The same machinery applies to RAG citations,
training masks, supervised target roles, DPO/RLHF pair consistency, and
evaluation harness prompts. A benchmark score should not be trusted if the
grader rubric, answer schema, system prompt, or private answer key can be
dropped or exposed by the interface.

\subsection{Provider Migration}

Provider migration is checked as a contract comparison, not as string matching.
PromptABI compares request and response surfaces: tool-call IDs, JSON-object
versus JSON-string arguments, parallel-call support, streaming chunk assembly,
stop semantics, context limits, structured-output modes, and error envelopes.
The output is a set of migration diagnostics with concrete fields and suggested
safe fixes.

\section{What the Tool Can Do}

Table~\ref{tab:capabilities} summarizes the main classes of checks. The table is
organized by programmer need rather than by implementation module.

\begin{table}[h]
\centering
\small
\begin{tabular}{p{0.24\linewidth}p{0.36\linewidth}p{0.28\linewidth}}
\toprule
Need & PromptABI check & Evidence produced \\
\midrule
Prevent prompt control injection &
Role-boundary non-forgeability over bounded template paths &
Forged marker, input field, rendered region, token evidence \\
\addlinespace
Avoid premature parser truncation &
Stop-overreachability over schemas, tools, and provider envelopes &
Stop firing point, parser state, malformed prefix \\
\addlinespace
Trust constrained decoding &
Tokenizer $\times$ grammar emptiness and ambiguity checks &
Token path, decoded text, grammar state, minimal string \\
\addlinespace
Migrate providers safely &
Provider and streaming contract comparison &
Lost field, changed encoding, response-shape witness \\
\addlinespace
Protect RAG, training, and eval invariants &
Context-budget, mask, role, and private-field obligations &
Survival proof, dropped-region witness, finite assignment \\
\addlinespace
Maintain reproducible CI &
Lockfiles, signed bundles, replayable solver certificates &
Artifact hash, diagnostic fingerprint, replay hash \\
\bottomrule
\end{tabular}
\caption{PromptABI capabilities stated as user-facing verification needs.}
\label{tab:capabilities}
\end{table}

\section{How It Proves Things}

PromptABI uses proof techniques that match the artifact being checked. It does
not use one universal solver for every problem.

\paragraph{Bounded symbolic execution.}
For supported chat-template fragments, PromptABI enumerates bounded control-flow
paths and loop counts. Each path is a sequence of symbolic segments annotated
with role ownership and source expressions. A finding is constructed only when a
user-controlled expression is in a region whose control markers are known.

\paragraph{Finite-language products.}
Tokenizers, grammars, and parser fragments are compiled to finite relations for
the witnesses under consideration. Products and projections answer emptiness,
ambiguity, and reachability questions with concrete paths.

\paragraph{SMT-backed finite contracts.}
Static contracts that combine budgets, roles, stops, masks, and provider facts
are lowered to finite SMT obligations. Solver models become counterexamples;
unsatisfiable obligations become proof summaries. Reduced solver-replay files
allow a result to be rerun without the original private artifacts.

\paragraph{Differential replay.}
Where a library is the specification, PromptABI runs small differential cases
against real tokenizer/template/parser behavior. This catches divergence between
the local abstraction and the implementation developers deploy.

\paragraph{Fail-closed abstention.}
If a construct is outside the supported fragment, a dependency is missing, a
schema cannot be modeled, or a source-derived value cannot be recovered,
PromptABI reports that fact. It does not silently turn an unsupported case into
a safe result.

\begin{theorem}[Bounded witness soundness]
For a supported checker fragment, every unsafe PromptABI diagnostic contains a
witness that is accepted by the artifact semantics used by that checker and
violates the stated interface obligation.
\end{theorem}

\begin{proof}[Proof sketch]
Each checker constructs diagnostics from semantic objects, not from ad-hoc
strings. The template checker builds a witness from an enumerated symbolic path,
role region, and input expression. The stop checker substitutes the stop into a
valid structured-output witness and records the parser state at truncation. The
grammar/tokenizer check records a concrete token path through the product. The
SMT layer records the satisfying assignment or unsat result for the lowered
finite obligation. Tests replay these witnesses through the same checker
semantics and, where applicable, real library adapters.
\end{proof}

\begin{theorem}[Honest fragment boundary]
PromptABI's successful safety claims are scoped to explicit supported fragments;
inputs outside those fragments produce abstentions or setup failures rather than
passing diagnostics.
\end{theorem}

\begin{proof}[Proof sketch]
Artifact loaders and checkers validate dialects, parser support, dependency
availability, schema fragments, and source-derived replay assumptions before
claiming success. Unsupported constructs are recorded in the diagnostic model,
and release tests require abstentions to remain visible.
\end{proof}

\section{Implementation}

PromptABI is a Python tool with a deterministic command-line interface and an
embedding API. The implementation is organized around typed artifacts and typed
diagnostics rather than around particular frameworks. The same verification
result can be rendered as text, JSON, SARIF, GitHub annotations, HTML reports,
signed bundles, or downstream integration payloads.

The CLI is intentionally CPU-only for the checks described here. That property
is important: users can run PromptABI in pull requests, air-gapped builds,
pre-commit hooks, release gates, fine-tuning preflight jobs, benchmark
publication workflows, and provider migration reviews without model weights or
provider credentials.

Privacy is part of the implementation contract. Diagnostics can expose exact
structural witnesses for public fixtures, but private workflows can choose
position-preserving redaction, hash-only evidence, or structural summaries while
preserving stable fingerprints.

\section{Evaluation}

PromptABI's evaluation separates verifier-maintenance tests from real-world
claims. Unit tests, property tests, and small local fixtures are useful for
maintaining the implementation, but they are not counted here as bug-finding
evidence. The evaluation in this paper counts only public upstream code that was
shipped in the wild and is pinned by repository, commit, path, license, and
SHA-256 hash.

\subsection{Research Questions}

The evaluation asks four questions.
\begin{enumerate}
  \item \textbf{Bug finding on real code:} when given exact production artifacts
  from public incidents, does PromptABI find the interface bug?
  \item \textbf{Attribution:} does the diagnostic identify the production
  artifact and boundary responsible for the bug?
  \item \textbf{Replayability:} can the result be rerun offline from pinned
  public source bytes without model weights, provider calls, or private data?
  \item \textbf{Cost:} is the check cheap enough for CI?
\end{enumerate}

\subsection{Corpus}

Table~\ref{tab:eval} lists the current evaluation corpus. Each row is exact code
in the wild. The first row is the upstream llama.cpp Phi chat-template file. The
second row is the upstream vLLM Qwen3 Coder tool-parser source. The tool loads
the pinned source excerpt, validates the source hash, extracts the template or
parser delimiters from those bytes, and only then runs the checker. If the
expected value cannot be recovered from the source, the replay fails.

\begin{table}[h]
\centering
\small
\begin{tabular}{p{0.22\linewidth}p{0.24\linewidth}p{0.25\linewidth}p{0.16\linewidth}}
\toprule
Production artifact & Public incident & Property checked & Result \\
\midrule
llama.cpp \texttt{microsoft-Phi-3.5-mini-instruct.jinja} at commit
\texttt{2016bf2} &
Phi-style role delimiter handling, llama.cpp PR 18462 &
Can raw message content forge role/special-token boundaries? &
Bug found: 18 role-boundary witnesses \\
\addlinespace
vLLM \texttt{qwen3coder\_tool\_parser.py} at commit \texttt{6a11d72} &
Qwen3 XML tool parser regressions, vLLM PR 40861 &
Can parser stop delimiters occur inside a valid tool argument boundary? &
Bug found: 2 stop-overreach witnesses \\
\bottomrule
\end{tabular}
\caption{Evaluation cases. Every case is replayed from exact pinned upstream
production code; synthetic reductions are not counted.}
\label{tab:eval}
\end{table}

\subsection{Results}

PromptABI finds both production-code bugs in the corpus. The Phi case yields 18
role-boundary witnesses after traversing 430 bytes of the exact upstream
template. The Qwen3 case yields 2 stop-overreach witnesses after extracting
\texttt{</tool\_call>}, \texttt{</function>}, and \texttt{</parameter>} from the
exact upstream parser source. On a local CPU run, the complete evaluation command
finishes in 0.64 seconds wall-clock time.

\begin{table}[h]
\centering
\small
\begin{tabular}{lr}
\toprule
Metric & Value \\
\midrule
Exact upstream production-code cases & 2 \\
Cases with expected bug found & 2 \\
Abstentions in counted cases & 0 \\
Model or provider calls & 0 \\
Total concrete witnesses & 20 \\
Wall-clock runtime & 0.64s \\
\bottomrule
\end{tabular}
\caption{Current production-code evaluation results.}
\label{tab:results}
\end{table}

\subsection{Interpretation}

The evaluation is intentionally small but strict. It does not claim that the
entire public real-bug corpus has been converted to production-code replay.
Non-production regression fixtures are excluded from the evaluation claim here.
The claim is: for the counted cases, PromptABI traverses the same public code
that shipped upstream, recovers the relevant boundary from that code, and
produces concrete diagnostics before inference.

\section{Limitations and Non-Goals}

PromptABI is not a model-behavior verifier. It does not prove that an LLM will
answer truthfully, refuse harmful requests, or choose a tool wisely. It proves
properties of the discrete interface that surrounds the model.

PromptABI is also intentionally bounded. Chat-template loops, grammar products,
schema recursion, SMT domains, and context-budget facts are finite. That makes
diagnostics replayable and CI-friendly, but it means claims must be read with
their bounds and supported-fragment metadata.

Finally, PromptABI is not a sandbox. It avoids provider calls and model
execution for verification, and it can operate on redacted evidence, but process
isolation and secret governance remain deployment responsibilities.

\section{Related Work}

PromptABI sits between several families of tools. Prompt linters and jailbreak
tests look for risky strings or behavioral failures, but they do not generally
compile tokenizer/template/provider artifacts into a shared semantics. Schema
validators and constrained decoders help one structured-output layer, but they
usually do not account for tokenizer drift, stop policies, provider envelopes,
or downstream parsers. Formal methods tools can prove finite obligations, but
they are rarely packaged around the artifacts LLM engineers already maintain.
PromptABI's contribution is to bring PL-style static and differential reasoning
to those concrete artifacts.

\section{Conclusion}

LLM systems need a verifier for the deterministic software boundary around the
model. PromptABI supplies that verifier. It treats prompt templates,
tokenizers, stops, schemas, tools, providers, budgets, training manifests, and
evaluation harnesses as a prompt ABI; compiles them to bounded semantics; and
emits replayable proofs, counterexamples, or abstentions before inference.

The resulting workflow is practical because it is CPU-only, version-aware,
source-mapped, privacy-conscious, and integrated with CI. It is credible because
its diagnostics carry executable witnesses and because its real-bug evaluation
traverses exact production code that shipped in the wild.

\end{document}
